\documentclass[10pt]{article}
% The assignment (Part I of the tutorial) as a standalone two-page handout.
% Keep it in sync with Part I of main.tex. No hyperref: references to the
% tutorial are written as section numbers (\S) of tutorial.pdf.

\usepackage[margin=1in]{geometry}
\usepackage{textcomp}
\usepackage{microtype}
\usepackage{amsmath,amssymb}
\usepackage{booktabs,tabularx,array}
\usepackage{enumitem}
\usepackage{xcolor}
\usepackage{fancyhdr}
\usepackage{lastpage}
\usepackage{xurl}
\usepackage[normalem]{ulem}

\definecolor{alertblue}{HTML}{1F4E79}
\newcommand{\alert}[1]{{\color{alertblue}\uline{#1}}}
\newcommand{\code}[1]{\path{#1}}
\newcommand{\MW}{\ensuremath{\,\mathrm{MW}}}
\newcommand{\MWh}{\ensuremath{\,\mathrm{MWh}}}
% Placeholder for a decision the instructor still has to fill in.
\newcommand{\instr}[1]{\textcolor{red}{\textbf{[Instructor: #1]}}}
\newcommand{\assignsection}[2]{\subsection*{#1}}

\pagestyle{fancy}
\setlength{\headheight}{15pt}
\fancyhf{}
\lhead{Data-Center Connection Planning Game}
\rhead{Assignment}
\cfoot{\thepage\ of \pageref{LastPage}}
\setlength{\parindent}{0pt}
\setlength{\parskip}{0.3em}
\setlist{itemsep=0pt,topsep=2pt}
\emergencystretch=3em

\begin{document}

\begin{center}
{\large ECE 51032: Computational Methods for Power System Analysis}\\[0.2em]
{\Large\bfseries Data-Center Connection Planning Game: Assignment}\\[0.2em]
Purdue University, Fall 2026
\end{center}

% This handout is the official assignment. The tutorial
% (\code{docs/student_tutorial/tutorial.pdf}) explains the game, the model, the
% tools, and strategy in detail.

% ---------------------------------------------------------------------------
\assignsection{The project}{asg:glance}

\paragraph{Your role.} Your team is a \alert{large-load developer}: you build
data centers, and your goal is to connect \alert{gigawatts of computing load}
to the Texas grid as quickly and as profitably as possible. The grid is a
2,000-bus synthetic model of Texas (\S7 of the tutorial), and the instructor plays the grid operator.

\paragraph{Interconnection options.} Each site is one row of a CSV file (\S8 of the tutorial):
\begin{itemize}[leftmargin=*]
  \item a \alert{bus} (where the site connects) and a \alert{connection
    capacity} $C$\MW{} (the load you ask the grid to serve);
  \item optionally, a \alert{co-located generator} of $G \le C$\MW{} that is
    used to offset load at the location, i.e., the net load for your load and
    generator there is $C - G$;
  \item optionally, a \alert{co-located 4-hour battery} of $P$\MW{} and
    $4P$\MWh, that will be operated by the system operator according to
    multiperiod economic dispatch to help integrate more load into the network
    (\S3 of the tutorial).
\end{itemize}
A team's sites may add up to at most \alert{1,000\MW} of connection capacity
per submission.

\paragraph{The game.} You are not alone on the grid. Every team's plan is
evaluated \alert{together} on the same network, to ensure the feasibility of
the grid during real time operation. If the plans do not all fit, the plans
will be curtailed down systematically to \alert{maximize the total load
interconnected} (\S2.2 of the tutorial). So the best place to build depends on
where the other teams build: a bus that looks empty may be crowded once
everyone picks it. After each submission you see every team's service fraction
and a grid map of the total load per bus, but never another team's plan
(\S14 of the tutorial).

% ---------------------------------------------------------------------------
\assignsection{Teams and communication}{asg:team}

Every team should have \alert{three students}, formed by yourselves. Each team
will be assigned a fixed \code{team_id} (chosen by yourselves) all game long,
and one member uploads each file. Questions, announcements, and finding
teammates go through the \alert{Discord} channel (link:
\url{https://discord.gg/Xwdgwvn5}). Teams may talk to and coordinate with each
other, but do not share your CSV or your exact plan with another team. You may
use AI tools (e.g., coding agents) for any part of the work and it is
encouraged to do so (\S5 of the tutorial). Your team is responsible for every number
it submits and must be able to explain its plan in the final presentation.

% ---------------------------------------------------------------------------
\assignsection{The three weeks}{asg:rules}

Every week plays the \alert{same game on the same data}: you are developing
data-center siting plans, possibly with co-located generation and storage,
that will be in effect for the next ten years, made in two submissions per
week (Fig.~1 of the tutorial). The practice week (Week 1) \alert{spans two
calendar weeks}, Nov 16--29, so that it does not overlap with Thanksgiving.
\begin{center}
\small
\begin{tabularx}{\linewidth}{@{}l l X X l@{}}
\toprule
Week & Dates & Submission 1 due & Submission 2 due & Grid energy cost\\
\midrule
1 (practice) & Nov 16--29 & Fri Nov 20, 11:59 PM & Fri Nov 27, 11:59 PM & none\\
2 (graded)   & Nov 30--Dec 6 & Thu Dec 3, 11:59 PM & Sun Dec 6, 11:59 PM & none\\
3 (graded)   & Dec 7--13 & Thu Dec 10, 11:59 PM & Sun Dec 13, 11:59 PM & at the LMPs\\
\bottomrule
\end{tabularx}
\end{center}
\begin{itemize}[leftmargin=*]
  \item \alert{Submission 1:} connect your first sites.
  \item \alert{Submission 2:} the Submission 1 results are released first, and
    every team's Submission 1 load/generation/storage becomes part of the base
    system. List \alert{only new sites} and they compete for the room that is
    left.
  \item Everything \alert{resets} between weeks, so you can use what you
    learned the week before.
  \item \alert{What you submit:} one CSV per team per submission, named
    \code{team{x}_week{x}_sub{x}.csv} (e.g., \code{team07_week2_sub1.csv}) and
    uploaded to \alert{BrightSpace}. Check it first with
    \code{ercot-game-check-csv} (\S8 and \S13 of the tutorial): an
    \alert{invalid or late file scores 0}. The free late days for the course
    cannot be used for the project.
  \item \alert{At the end:} a short team presentation (last week of class) and
    a report on what you did, why, and optionally, some insights on explaining
    the results of the game.
\end{itemize}

% ---------------------------------------------------------------------------
\assignsection{Scoring rules}{asg:scoring}

Each submission is scored in US dollars, once, for the ten-year value of the
sites it builds (\S2 of the tutorial; worked example in \S2.1):
score $=$ IT value $-$ energy cost $-$ generator operating cost $-$ generator
capital $-$ storage capital $-$ site cost.
\begin{center}
\small
\begin{tabularx}{\linewidth}{@{}l l X@{}}
\toprule
Term & Value & How it is charged\\
\midrule
IT value & \$1,000 per MWh & For every MWh of your load \alert{that is actually served}.\\
Energy cost & LMP at your site & Grid purchases times the LMP. \alert{Week 3 only.}\\
Generator operation & \$80/\$150/\$250 per MWh & First 25\% / next 50\% / last 25\% of capacity.\\
Generator capital & \$2.0M per MW & Once, on the allocated generator capacity $\alpha G$.\\
Storage capital & \$400k/MW + \$250k/MWh & Once, on $\alpha P$ and $\alpha B$; a 4-hour battery costs \$1.4M per MW.\\
Site cost & \$25M per MW & Once, on the allocated connection capacity $\alpha C$.\\
\bottomrule
\end{tabularx}
\end{center}
Here $\alpha \in [0,1]$ is a site's allocated share. The representative days make up one typical year, and the ten
years count in \alert{net present value} at a 7\% discount rate:
$\text{NPV} = \sum_{y=1}^{10} \text{(yearly operating net revenue)}/1.07^{y}
- \text{capital}$, where $\sum_{y=1}^{10} 1.07^{-y} = 7.02$. So one MW served in
every hour earns $8{,}760 \times \$1{,}000 \times 7.02 = \alert{\$61.5\text{M}}$,
against \$25M of site cost. Capital is paid up front and is not discounted.

% ---------------------------------------------------------------------------
\assignsection{Grading}{asg:conduct}

The assignment counts for \alert{40 points} of the final grade:
\begin{itemize}[leftmargin=*]
  \item \alert{30 points}: the presentation, the report, and your work.
  \item \alert{10 points}: your final game score $S$, the sum of your four
    Week 2 and Week 3 submission scores.
    \begin{itemize}
      \item \alert{5 absolute points}: $5 \times \min\!\big(1, \max(0,
        S/(2B))\big)$, where $B = \$49.67$B is the final score $S$ of a
        reference team that runs 1,000\MW{} off the grid ($G = C$): \$24.84B in
        each of W2S1 and W3S1, and \$0 in W2S2 and W3S2 (no new sites).
      \item \alert{5 rank points}: 5, 4, 3, 2, or 1 for the top 20\%,
        20--40\%, 40--60\%, 60--80\%, or bottom 20\% of teams.
    \end{itemize}
  \item \alert{Bonus, up to 5 points}: 3, 2, and 1 for the top three final
    scores, and up to 2 for reporting non-trivial and non-intended
    issues/bugs in the project design.
\end{itemize}
Purdue's academic integrity policy applies. To question an evaluation, email
Minghao Mou (\code{mmou@purdue.edu}) within 7 days of the results release.

% ---------------------------------------------------------------------------
\assignsection{Resources and links}{asg:resources}

\begin{itemize}[leftmargin=*]
  \item \alert{Tutorial} (\code{docs/student_tutorial/tutorial.pdf} in the repository): setup (\S4), an AI-agent walkthrough (\S6), and strategy (\S11).
  \item \alert{Code and data}:
    \url{https://github.com/MhaoMou/ercot-connection-planning-instructor}.
  \item \alert{Discord}: \url{https://discord.gg/Xwdgwvn5}.
\end{itemize}

\end{document}
